% TOP_PROBLEM: {"id": 3075, "title": "Average diameter of a bounded cell of a simple arrangement", "category_id": 6, "category": {"id": 6, "name": "geometry", "display_name": "Geometry", "description": "Euclidean and non-Euclidean geometry, geometric structures.", "slug": "geometry", "order_index": 6, "created_at": "2026-07-31T15:26:25.671Z"}, "status": "open", "problem_number": "OPG-605", "set_id": 12}
% FIRST_POSTED: 2026-10-01
% ENTRY_KIND: statement_only
% UnsolvedMath ID 3075; source catalogue code OPG-605
% !TeX program = lualatex
% Definition and sources only; this file is not an LLM solution attempt.
\documentclass[11pt,a4paper]{article}
\usepackage[margin=25mm]{geometry}
\usepackage{fontspec}
\setmainfont{lmroman10-regular.otf}[BoldFont=lmroman10-bold.otf,
 ItalicFont=lmroman10-italic.otf,BoldItalicFont=lmroman10-bolditalic.otf]
\usepackage{amsmath,amssymb,mathtools}
\usepackage{unicode-math}
\setmathfont{latinmodern-math.otf}
\AtBeginDocument{\let\setminus\smallsetminus}
\usepackage{xurl,hyperref}
\hypersetup{colorlinks=true,urlcolor=blue}
\setcounter{secnumdepth}{0}
\setlength{\parindent}{0pt}
\setlength{\parskip}{5pt}
\setlength{\emergencystretch}{3em}
\begin{document}
\section*{Average diameter of a bounded cell of a simple arrangement}\label{3075}
{\small UnsolvedMath ID 3075 · OPG-605 · Geometry · OpenGarden}
\subsection{Definitions and mathematical statement}
Fix integers $d\ge1$ and $n\ge d+1$. A simple affine hyperplane
arrangement in $\mathbb R^d$ is a family $\mathcal A$ of $n$ hyperplanes
in general position: every $j\le d$ distinct hyperplanes meet in an
affine subspace of dimension $d-j$, and no $d+1$ have a common point.
The components of their complement are open polyhedral regions.
The closures of the bounded regions are full-dimensional convex polytopes,
called the bounded cells.

For such a cell $P$, let $\delta(P)$ be the diameter of its edge graph:
the maximum, over pairs of vertices, of the smallest number of polytope
edges in a path joining them. This is a combinatorial graph distance,
not a Euclidean distance between points. List all bounded cells as
$P_1,\ldots,P_I$ and define
\[
 \overline\delta(\mathcal A)=\frac1I\sum_{j=1}^I\delta(P_j).
\]
The general-position assumptions give $I=\binom{n-1}{d}>0$, so the
average is defined and weights every bounded cell equally.

The conjecture states that
$\overline\delta(\mathcal A)\le d$ for every such arrangement, independently
of $n$. Individual cells are allowed to have diameter greater than $d$;
the requirement bounds only the mean over all bounded cells.
Unbounded regions are omitted, and neither cell volume nor number of
vertices changes its weight in the average. Establishing a bound for just
one selected polytope therefore does not settle the arrangement statement.
\subsection{Short English statement}
Is the average graph diameter of all bounded cells of a simple hyperplane arrangement at most the ambient dimension?
\subsection{Sources}
\begin{itemize}
\item[S1] ULAM.AI and the UnsolvedMath Contributors, supplied problems.json, record 3075 (OPG-605), OpenGarden collection. Catalogue identity and source transcription; CC BY 4.0.
\url{https://huggingface.co/datasets/ulamai/UnsolvedMath}.
\item[S2] Open Problem Garden, Average diameter of a bounded cell of a simple arrangement. Original problem and discussion; the mathematical formulation below is expanded from the supplied source record.
\url{http://www.openproblemgarden.org/op/average_diameter_of_a_bounded_cell_of_a_simple_arrangement}.
\end{itemize}
\end{document}
